\theory{Ring theory}{What happens when addition and multiplication coexist in one algebraic system, even when multiplication is not reversible or commutative?}{Ring theory studies systems modeled on the integers. It studies ideals, homomorphisms, modules, polynomial identities, commutative and noncommutative rings, and the structure of special rings such as matrix rings and division rings.}{Algebraic geometry, algebraic number theory, coding theory, invariant theory, noncommutative geometry, linear algebra, and homological algebra.}{Addition, multiplication, ideals, and modules describe algebraic systems in which factorization and quotient structure can be studied together.}

\paragraph{The idea and history}
Integers can be added, subtracted, and multiplied. A ring keeps these operations but does not require every nonzero element to have a multiplicative inverse. Polynomial rings, matrices, functions, modular arithmetic, and coordinate rings are all examples. The theory developed from algebraic number theory, algebraic geometry, invariant theory, and attempts to generalize complex numbers \cite{ring_ref1}.

Hamilton's quaternions and other hypercomplex systems led mathematicians to matrix and noncommutative examples. Joseph Wedderburn identified finite-dimensional algebras with matrix structures, and Emil Artin generalized this work to Artinian rings. Emmy Noether's work on ideals and ascending chain conditions led to the term Noetherian ring and transformed modern algebra \cite{ring_ref2}.

\end{historylist}
\section*{3. Theory}
\subsection*{Core theory}
\paragraph{Rings, ideals, and quotients}
A ring $R$ has an abelian group structure under addition and an associative multiplication distributing over addition. It may have a multiplicative identity. A commutative ring satisfies $ab=ba$; a noncommutative ring need not.

An ideal $I$ is an additive subgroup absorbing multiplication: $r i$ and $i r$ belong to $I$ whenever $i\in I$ and $r\in R$. Ideals allow a quotient ring $R/I$, in which elements differing by something in $I$ are identified. A ring homomorphism has a kernel that is an ideal, and the first isomorphism theorem says the image is isomorphic to the quotient by that kernel \cite{ring_ref1}.

Prime and maximal ideals generalize prime numbers and fields. In a commutative ring, a prime ideal $P$ satisfies $ab\in P$ only when $a\in P$ or $b\in P$. A maximal ideal has a field as quotient. These notions are the algebraic basis of the spectrum of a ring in algebraic geometry.

\paragraph{Modules and important classes}
A module over $R$ is like a vector space in which scalars come from a ring instead of a field. Modules are the natural language for representations of rings and for systems of linear equations over the integers. A ring is Noetherian when every ascending chain of ideals stops, or equivalently when every ideal is finitely generated.

An integral domain is a commutative ring without zero divisors. A field is a domain in which every nonzero element is invertible. A principal ideal domain has every ideal generated by one element, while a Euclidean domain has a division algorithm. The integers form a Euclidean domain; polynomial rings over fields often do as well.

Noncommutative rings include square matrix rings, endomorphism rings, group rings, and universal enveloping algebras. Their modules can be more complicated because left and right ideals differ. Morita theory studies when two rings have equivalent module categories, showing that the category of modules can be more important than the ring's particular presentation \cite{ring_ref3}.

\paragraph{Structure and applications}
The Artin--Wedderburn theorem describes semisimple rings as products of matrix rings over division rings. Wedderburn's little theorem says every finite division ring is commutative, so it is a field. Nakayama's lemma controls generators of finitely generated modules over local rings. These results let algebraists replace a complicated ring by standard building blocks.

For an affine algebraic variety, regular functions form a coordinate ring. The geometry of the variety can be translated into ideal and quotient questions. The Hilbert Nullstellensatz connects maximal ideals of a polynomial ring over an algebraically closed field with points of the variety. Symmetric polynomials form an invariant ring; the fundamental theorem says they are generated by the elementary symmetric polynomials \cite{ring_ref4}.

\paragraph{Application: Reed--Solomon codes}
A concrete connection between ring theory and coding theory arises through Reed--Solomon codes. Let $\mathbb F_q$ be a finite field with $q$ elements, fix distinct elements $\alpha_1,\ldots,\alpha_n\in\mathbb F_q$ with $n\le q-1$, and choose an integer $k\le n$. The ring of polynomials $\mathbb F_q[x]$ provides the algebraic setting. A message is a polynomial $m(x)\in\mathbb F_q[x]$ of degree less than $k$. The encoder evaluates $m$ at the fixed points, producing the codeword
\[
c=(m(\alpha_1),m(\alpha_2),\ldots,m(\alpha_n)).
\]
The set of all such codewords forms a linear code of dimension $k$ and length $n$. Two distinct degree-$(k{-}1)$ polynomials agree on at most $k{-}1$ points, so any two codewords differ in at least $n{-}k{+}1$ positions. This minimum distance guarantees that the code can correct up to $\lfloor(n{-}k)/2\rfloor$ symbol errors.

The algebraic structure is essential: the set of codewords is an ideal in the ring $\mathbb F_q^n$ (under componentwise operations), and decoding can be performed by solving a system of polynomial congruences in $\mathbb F_q[x]$. Euclidean division in the polynomial ring provides the key algorithmic step. Reed--Solomon codes are used in QR codes, deep-space communication, and data storage systems such as CDs and DVDs \cite{ring_ref1}.

\section*{4. Important theorems and results}
\begin{enumerate}[leftmargin=*,itemsep=0.35em]
\item \textbf{Ring isomorphism theorems.} Kernels, images, quotients, and induced maps obey analogues of the group isomorphism theorems.

\item \textbf{Nakayama's lemma.} Over a local ring, generators of a finitely generated module can be detected after reducing modulo the maximal ideal.

\item \textbf{Artin--Wedderburn theorem.} A semisimple Artinian ring is a finite product of full matrix rings over division rings.

\item \textbf{Wedderburn's little theorem.} Every finite division ring is a field.

\item \textbf{Hilbert basis theorem.} If $R$ is Noetherian, then the polynomial ring $R[x]$ is Noetherian.

\item \textbf{Hilbert Nullstellensatz.} Over an algebraically closed field, maximal ideals of a polynomial ring correspond to points, and ideals describe the polynomial equations vanishing on sets.
\end{enumerate}
\noindent\emph{Source for these results:} \cite{ring_ref1}.

\theoryapplications
\theoryconnections{ring_theory}{%
\item Groups describe addition and multiplication separately, while ring theory studies when these operations interact through distributivity.
\item Modules generalize vector spaces when coefficients come from a ring rather than a field, and ideals identify which expressions should count as zero in a quotient.
\item Polynomial rings turn equations into ring elements, so factoring, divisibility, and quotienting become tools for studying geometric solution sets \cite{ring_ref1,ring_ref2}.
}{%
\item Ring theory helps algebraic geometry by translating a variety into its coordinate ring, where geometric inclusion becomes ideal inclusion.
\item In number theory, ideals restore a workable factorization theory in rings where elements do not factor uniquely.
\item Modules support representations and homological algebra, while quotient rings over finite fields provide the arithmetic used to construct and decode error-correcting codes \cite{ring_ref1,ring_ref3}.
}
\section*{7. Sample Questions}
\emph{Exercise reference:} \cite{ring_ref1}. The cited edition is the source for the exercise set; consult its Exercises section for the official statement and numbering.
\begin{enumerate}[leftmargin=*,itemsep=0.9em]
\item \textbf{Exercise 1. Find the ideals of $\mathbb Z$.} Every ideal is $n\mathbb Z$ for a unique nonnegative integer $n$. If an ideal contains a least positive element $n$, division with remainder shows every element is a multiple of $n$.

\item \textbf{Exercise 2. Describe $\mathbb Z/6\mathbb Z$.} Its elements are the six residue classes. The class of $2$ is a zero divisor because $2\cdot3=0$ modulo $6$, so the quotient is not a field.

\item \textbf{Exercise 3. Why is $2\mathbb Z$ an ideal of $\mathbb Z$?} It is closed under addition and negatives, and multiplying an even number by any integer remains even.

\item \textbf{Exercise 4. Is $\mathbb Z/5\mathbb Z$ a field?} Yes. Every nonzero class has an inverse because $5$ is prime; for example $2^{-1}=3$ since $2\cdot3=6\equiv1$ modulo $5$.

\item \textbf{Exercise 5. Why are matrices a noncommutative ring?} Matrix addition and multiplication satisfy the ring laws, but usually $AB\ne BA$. The failure of commutativity is exactly why left and right ideals and modules must be distinguished.

\end{enumerate}
\section*{8. References}
\begin{thebibliography}{9}
\bibitem{ring_ref1} T. Y. Lam, \emph{A First Course in Noncommutative Rings}, 2nd ed., Springer, 2001.
\bibitem{ring_ref2} C. Faith, \emph{Rings and Things and a Fine Array of Twentieth-Century Associative Algebra}, AMS, 2004.
\bibitem{ring_ref3} T. Y. Lam, \emph{Lectures on Modules and Rings}, Springer, 1999.
\bibitem{ring_ref4} D. Eisenbud, \emph{Commutative Algebra with a View Toward Algebraic Geometry}, Springer, 1995.
\end{thebibliography}
